\documentclass[10pt,a4paper]{amsart}
\usepackage{amsmath,amssymb,amsthm,graphicx}
% For dealing with references we use the comment environment
\usepackage{verbatim}
\newenvironment{reference}{\comment}{\endcomment}
%\newenvironment{reference}{}{}
\newenvironment{slogan}{\comment}{\endcomment}
\newenvironment{history}{\comment}{\endcomment}
% For commutative diagrams we use Xy-pic
\usepackage[all]{xy}

% We use 2cell for 2-commutative diagrams.
\xyoption{2cell}
\UseAllTwocells
\theoremstyle{plain}
\newtheorem{theorem}[subsection]{Theorem}
\newtheorem{proposition}[subsection]{Proposition}
\newtheorem{lemma}[subsection]{Lemma}

\theoremstyle{definition}
\newtheorem{definition}[subsection]{Definition}
\newtheorem{example}[subsection]{Example}
\newtheorem{exercise}[subsection]{Exercise}
\newtheorem{situation}[subsection]{Situation}

\theoremstyle{remark}
\newtheorem{remark}[subsection]{Remark}
\newtheorem{remarks}[subsection]{Remarks}

\numberwithin{equation}{subsection}
% Macros
%
\def\lim{\mathop{\mathrm{lim}}\nolimits}
\def\colim{\mathop{\mathrm{colim}}\nolimits}
\def\Spec{\mathop{\mathrm{Spec}}}
\def\Hom{\mathop{\mathrm{Hom}}\nolimits}
\def\Ext{\mathop{\mathrm{Ext}}\nolimits}
\def\SheafHom{\mathop{\mathcal{H}\!\mathit{om}}\nolimits}
\def\SheafExt{\mathop{\mathcal{E}\!\mathit{xt}}\nolimits}
\def\Sch{\mathit{Sch}}
\def\Mor{\mathop{\mathrm{Mor}}\nolimits}
\def\Ob{\mathop{\mathrm{Ob}}\nolimits}
\def\Sh{\mathop{\mathit{Sh}}\nolimits}
\def\NL{\mathop{N\!L}\nolimits}
\def\CH{\mathop{\mathrm{CH}}\nolimits}
\def\proetale{{pro\text{-}\acute{e}tale}}
\def\etale{{\acute{e}tale}}
\def\QCoh{\mathit{QCoh}}
\def\Ker{\mathop{\mathrm{Ker}}}
\def\Im{\mathop{\mathrm{Im}}}
\def\Coker{\mathop{\mathrm{Coker}}}
\def\Coim{\mathop{\mathrm{Coim}}}

% Boxtimes
%
\DeclareMathSymbol{\boxtimes}{\mathbin}{AMSa}{"02}

%
% Macros for moduli stacks/spaces
%
\def\QCohstack{\mathcal{QC}\!\mathit{oh}}
\def\Cohstack{\mathcal{C}\!\mathit{oh}}
\def\Spacesstack{\mathcal{S}\!\mathit{paces}}
\def\Quotfunctor{\mathrm{Quot}}
\def\Hilbfunctor{\mathrm{Hilb}}
\def\Curvesstack{\mathcal{C}\!\mathit{urves}}
\def\Polarizedstack{\mathcal{P}\!\mathit{olarized}}
\def\Complexesstack{\mathcal{C}\!\mathit{omplexes}}
% \Pic is the operator that assigns to X its picard group, usage \Pic(X)
% \Picardstack_{X/B} denotes the Picard stack of X over B
% \Picardfunctor_{X/B} denotes the Picard functor of X over B
\def\Pic{\mathop{\mathrm{Pic}}\nolimits}
\def\Picardstack{\mathcal{P}\!\mathit{ic}}
\def\Picardfunctor{\mathrm{Pic}}
\def\Deformationcategory{\mathcal{D}\!\mathit{ef}}
\usepackage{cleveref}
\crefname{originallemma}{original lemma}{original lemmas}
\Crefname{originallemma}{Original lemma}{Original lemmas}
\makeatletter
\newlabel{stacks-sheaves-lemma-QC-compare-fppf}{{0H0X}{source}{}{sourcepointer-1}{}}
\newlabel{stacks-sheaves-lemma-QC-compare-fppf@cref}{{[originallemma][0][]0H0X}{[1][0][]source}{}{}{}}
\makeatother
\crefname{originalequation}{original equation}{original equations}
\Crefname{originalequation}{Original equation}{Original equations}
\makeatletter
\newlabel{stacks-sheaves-equation-alternative-ringed}{{0H0C}{source}{}{sourcepointer-2}{}}
\newlabel{stacks-sheaves-equation-alternative-ringed@cref}{{[originalequation][0][]0H0C}{[1][0][]source}{}{}{}}
\makeatother
\crefname{originallemma}{original lemma}{original lemmas}
\Crefname{originallemma}{Original lemma}{Original lemmas}
\makeatletter
\newlabel{lemma-bousfield-colocalization}{{07BA}{source}{}{sourcepointer-3}{}}
\newlabel{lemma-bousfield-colocalization@cref}{{[originallemma][0][]07BA}{[1][0][]source}{}{}{}}
\makeatother
\crefname{originalremark}{original remark}{original remarks}
\Crefname{originalremark}{Original remark}{Original remarks}
\makeatletter
\newlabel{remark-QCoh-admissible}{{0H11}{source}{}{sourcepointer-4}{}}
\newlabel{remark-QCoh-admissible@cref}{{[originalremark][0][]0H11}{[1][0][]source}{}{}{}}
\makeatother
\crefname{originallemma}{original lemma}{original lemmas}
\Crefname{originallemma}{Original lemma}{Original lemmas}
\makeatletter
\newlabel{sites-cohomology-lemma-cartesion-plus-topology}{{0GZS}{source}{}{sourcepointer-5}{}}
\newlabel{sites-cohomology-lemma-cartesion-plus-topology@cref}{{[originallemma][0][]0GZS}{[1][0][]source}{}{}{}}
\makeatother
\crefname{originallemma}{original lemma}{original lemmas}
\Crefname{originallemma}{Original lemma}{Original lemmas}
\makeatletter
\newlabel{stacks-sheaves-lemma-quasi-coherent-alternative}{{0H0E}{source}{}{sourcepointer-6}{}}
\newlabel{stacks-sheaves-lemma-quasi-coherent-alternative@cref}{{[originallemma][0][]0H0E}{[1][0][]source}{}{}{}}
\makeatother
\crefname{originallemma}{original lemma}{original lemmas}
\Crefname{originallemma}{Original lemma}{Original lemmas}
\makeatletter
\newlabel{sites-cohomology-lemma-section-RHom-over-U}{{08JA}{source}{}{sourcepointer-7}{}}
\newlabel{sites-cohomology-lemma-section-RHom-over-U@cref}{{[originallemma][0][]08JA}{[1][0][]source}{}{}{}}
\makeatother
\crefname{originallemma}{original lemma}{original lemmas}
\Crefname{originallemma}{Original lemma}{Original lemmas}
\makeatletter
\newlabel{sites-cohomology-lemma-QC-hom-out-of-plus-topology}{{0H0T}{source}{}{sourcepointer-8}{}}
\newlabel{sites-cohomology-lemma-QC-hom-out-of-plus-topology@cref}{{[originallemma][0][]0H0T}{[1][0][]source}{}{}{}}
\makeatother
\crefname{originallemma}{original lemma}{original lemmas}
\Crefname{originallemma}{Original lemma}{Original lemmas}
\makeatletter
\newlabel{lemma-cohomology-parasitic}{{0H13}{source}{}{sourcepointer-9}{}}
\newlabel{lemma-cohomology-parasitic@cref}{{[originallemma][0][]0H13}{[1][0][]source}{}{}{}}
\makeatother
\crefname{originallemma}{original lemma}{original lemmas}
\Crefname{originallemma}{Original lemma}{Original lemmas}
\makeatletter
\newlabel{lemma-cohomology-parasitic-complex}{{0H14}{source}{}{sourcepointer-10}{}}
\newlabel{lemma-cohomology-parasitic-complex@cref}{{[originallemma][0][]0H14}{[1][0][]source}{}{}{}}
\makeatother
\crefname{originallemma}{original lemma}{original lemmas}
\Crefname{originallemma}{Original lemma}{Original lemmas}
\makeatletter
\newlabel{lemma-higher-shriek-QC}{{0H15}{source}{}{sourcepointer-11}{}}
\newlabel{lemma-higher-shriek-QC@cref}{{[originallemma][0][]0H15}{[1][0][]source}{}{}{}}
\makeatother
\title{Proposition qc compare in stacks-perfect.tex}
\author{The Stacks Project Authors}
\date{stacks-project:a04446e57ec1fbc252a871afcec7752fb2807b14; GNU Free Documentation License 1.2 (as published in the Stacks Project COPYING)}
\begin{document}
\maketitle
\noindent\textit{Editorial note.} All mathematical statement and proof text is reproduced verbatim from the author's chapter stacks-perfect.tex of the Stacks Project, fixed at commit a04446e57ec1fbc252a871afcec7752fb2807b14. The pinned preamble is supplied as its own file; the wrapper class, the xr-hyper wiring and the \textbackslash{}externaldocument list are omitted because the assembler supplies the class.\par\medskip
\noindent\textit{Editorial note.} Numbering inside this standalone excerpt is editorial, not the full Stacks book print numbering.\par\medskip
\noindent\textit{Editorial note.} The candidate is from the corrected band-I follow-on screen (line C) of the extractor's stacks-candidates line: its author label is absent from the live master in every form searched and it belongs to none of the earlier judged pools. The H2 judgement and its reason were made in this module by reading the author's own statement and complete proof in the fixed checkout. Original candidate key: stacks:stacks-perfect.tex:proposition-QC-compare.\par\medskip
\noindent\textit{Editorial note.} Reference labels in the proof resolve to navigation pointers to the original author source; these pointers do not reproduce proofs of the prerequisites.\par\medskip
\noindent\textit{Allowed standard prerequisites (editorial).}\begin{itemize}
\item Prerequisites named by the author inside the counted span. The counted span is reproduced verbatim, including the author's own citations. Its prerequisites are the earlier results of the same fixed Stacks Project version that the author cites there and in the chapter's lead-in.
\end{itemize}\medskip
\ifdefined\hypertarget\hypertarget{sourcepointer-1}{}\fi
\noindent\textit{Source reference.} Original lemma 0H0X: sources/stacks-sheaves.tex, source lines 4964--4964. Navigation pointer used in the counted span; the span contains the pointer only and the cited result is not reproduced inside it. Registered as a non-load-bearing navigation stub in the established form for Stacks material. No bibliography citation is involved.\par\medskip
\ifdefined\hypertarget\hypertarget{sourcepointer-2}{}\fi
\noindent\textit{Source reference.} Original equation 0H0C: sources/stacks-sheaves.tex, source lines 4510--4510. Navigation pointer used in the counted span; the span contains the pointer only and the cited result is not reproduced inside it. Registered as a non-load-bearing navigation stub in the established form for Stacks material. No bibliography citation is involved.\par\medskip
\ifdefined\hypertarget\hypertarget{sourcepointer-3}{}\fi
\noindent\textit{Source reference.} Original lemma 07BA: sources/stacks-perfect.tex, source lines 862--862. Navigation pointer used in the counted span; the span contains the pointer only and the cited result is not reproduced inside it. Registered as a non-load-bearing navigation stub in the established form for Stacks material. No bibliography citation is involved.\par\medskip
\ifdefined\hypertarget\hypertarget{sourcepointer-4}{}\fi
\noindent\textit{Source reference.} Original remark 0H11: sources/stacks-perfect.tex, source lines 909--909. Navigation pointer used in the counted span; the span contains the pointer only and the cited result is not reproduced inside it. Registered as a non-load-bearing navigation stub in the established form for Stacks material. No bibliography citation is involved.\par\medskip
\ifdefined\hypertarget\hypertarget{sourcepointer-5}{}\fi
\noindent\textit{Source reference.} Original lemma 0GZS: sources/sites-cohomology.tex, source lines 11608--11608. Navigation pointer used in the counted span; the span contains the pointer only and the cited result is not reproduced inside it. Registered as a non-load-bearing navigation stub in the established form for Stacks material. No bibliography citation is involved.\par\medskip
\ifdefined\hypertarget\hypertarget{sourcepointer-6}{}\fi
\noindent\textit{Source reference.} Original lemma 0H0E: sources/stacks-sheaves.tex, source lines 4534--4534. Navigation pointer used in the counted span; the span contains the pointer only and the cited result is not reproduced inside it. Registered as a non-load-bearing navigation stub in the established form for Stacks material. No bibliography citation is involved.\par\medskip
\ifdefined\hypertarget\hypertarget{sourcepointer-7}{}\fi
\noindent\textit{Source reference.} Original lemma 08JA: sources/sites-cohomology.tex, source lines 8874--8874. Navigation pointer used in the counted span; the span contains the pointer only and the cited result is not reproduced inside it. Registered as a non-load-bearing navigation stub in the established form for Stacks material. No bibliography citation is involved.\par\medskip
\ifdefined\hypertarget\hypertarget{sourcepointer-8}{}\fi
\noindent\textit{Source reference.} Original lemma 0H0T: sources/sites-cohomology.tex, source lines 11678--11678. Navigation pointer used in the counted span; the span contains the pointer only and the cited result is not reproduced inside it. Registered as a non-load-bearing navigation stub in the established form for Stacks material. No bibliography citation is involved.\par\medskip
\ifdefined\hypertarget\hypertarget{sourcepointer-9}{}\fi
\noindent\textit{Source reference.} Original lemma 0H13: sources/stacks-perfect.tex, source lines 1154--1154. Navigation pointer used in the counted span; the span contains the pointer only and the cited result is not reproduced inside it. Registered as a non-load-bearing navigation stub in the established form for Stacks material. No bibliography citation is involved.\par\medskip
\ifdefined\hypertarget\hypertarget{sourcepointer-10}{}\fi
\noindent\textit{Source reference.} Original lemma 0H14: sources/stacks-perfect.tex, source lines 1176--1176. Navigation pointer used in the counted span; the span contains the pointer only and the cited result is not reproduced inside it. Registered as a non-load-bearing navigation stub in the established form for Stacks material. No bibliography citation is involved.\par\medskip
\ifdefined\hypertarget\hypertarget{sourcepointer-11}{}\fi
\noindent\textit{Source reference.} Original lemma 0H15: sources/stacks-perfect.tex, source lines 1197--1197. Navigation pointer used in the counted span; the span contains the pointer only and the cited result is not reproduced inside it. Registered as a non-load-bearing navigation stub in the established form for Stacks material. No bibliography citation is involved.\par\medskip
\begin{proposition}
\label{proposition-QC-compare}
Let $\mathcal{X}$ be an algebraic stack. Then $\mathit{QC}(\mathcal{X})$
is canonically equivalent to $D_\QCoh(\mathcal{O}_\mathcal{X})$.
\end{proposition}

\begin{proof}
By Sheaves on Stacks, Lemma \ref{stacks-sheaves-lemma-QC-compare-fppf}
pullback by the comparison morphism
$\epsilon : \mathcal{X}_{affine, fppf} \to \mathcal{X}_{affine}$
identifies $\mathit{QC}(\mathcal{X})$ with a full subcategory
$Q_\mathcal{X} \subset D(\mathcal{X}_{affine, fppf}, \mathcal{O})$.
Using the equivalence of ringed topoi in
Sheaves on Stacks, Equation (\ref{stacks-sheaves-equation-alternative-ringed})
we may and do view $Q_\mathcal{X}$ as a full subcategory of
$D(\mathcal{X}_{fppf}, \mathcal{O})$.

\medskip\noindent
Similarly by Lemma \ref{lemma-bousfield-colocalization} and
Remark \ref{remark-QCoh-admissible} we find that
$D_\QCoh(\mathcal{O}_\mathcal{X})$ may be viewed as the
left orthogonal $\mathcal{A}$ of the left admissible subcategory
$D_{\textit{Parasitic} \cap \textit{LQCoh}^{fbc}}(\mathcal{O}_\mathcal{X})$
of $D_{\textit{LQCoh}^{fbc}}(\mathcal{O}_\mathcal{X})$.

\medskip\noindent
To finish we will show that $Q_\mathcal{X}$ is equal to
$\mathcal{A}$ as subcategories of $D(\mathcal{X}_{fppf}, \mathcal{O})$.

\medskip\noindent
Step 1: $Q_\mathcal{X}$ is contained in
$D_{\textit{LQCoh}^{fbc}}(\mathcal{O}_\mathcal{X})$.
An object $K$ of $Q_\mathcal{X}$ is characterized by the property
that $K$, viewed as an object of $D(\mathcal{X}_{affine, fppf}, \mathcal{O})$
satisfies $R\epsilon_*K$ is an object of
$\mathit{QC}(\mathcal{X}_{affine}, \mathcal{O})$.
This in turn means exactly that for all morphisms $x \to x'$
of $\mathcal{X}_{affine}$ the map
$$
R\Gamma(x', K) \otimes_{\mathcal{O}(x')}^\mathbf{L} \mathcal{O}(x)
\longrightarrow
R\Gamma(x, K)
$$
is an isomorphism, see footnote in statement of Cohomology on Sites,
Lemma \ref{sites-cohomology-lemma-cartesion-plus-topology}.
Now, if $x' \to x$ lies over a flat morphism of affine schemes,
then this means that
$$
H^i(x', K) \otimes_{\mathcal{O}(x')} \mathcal{O}(x)
\cong
H^i(x, K)
$$
This clearly means that $H^i(K)$ is a sheaf for the \'etale topology
(Sheaves on Stacks, Lemma \ref{stacks-sheaves-lemma-quasi-coherent-alternative})
and that it has the flat base change property (small detail omitted).

\medskip\noindent
Step 2: $Q_\mathcal{X}$ is contained in $\mathcal{A}$.
To see this it suffices to show that for $K$ in $Q_\mathcal{X}$
we have $\Hom(K, P) = 0$ for all $P$ in
$D_{\textit{Parasitic} \cap \textit{LQCoh}^{fbc}}(\mathcal{O}_\mathcal{X})$.
Consider the object
$$
H = R\SheafHom_{\mathcal{O}_\mathcal{X}}(K, P)
$$
Let $x$ be an object of $\mathcal{X}$ which lies over an affine scheme
$U = p(x)$. By Cohomology on Sites, Lemma
\ref{sites-cohomology-lemma-section-RHom-over-U} we have
the first equality in
$$
R\Gamma(x, H) =
R\Hom_{\mathcal{O}_\mathcal{X}}(K|_{\mathcal{X}/x}, P|_{\mathcal{X}/x}) =
R\Hom_{\mathcal{O}}(K|_{\mathcal{X}_{affine}/x}, P|_{\mathcal{X}_{affine}/x})
$$
The second equality stems from the fact that the topos of the site
$\mathcal{X}/x$ is equivalent to the topos of the site
$\mathcal{X}_{affine}/x$, see Sheaves on Stacks, Equation
(\ref{stacks-sheaves-equation-alternative-ringed}).
We may write $K = \epsilon^*N$ for some $N$ in $\mathit{QC}(\mathcal{O})$.
Then by Cohomology on Sites, Lemma
\ref{sites-cohomology-lemma-QC-hom-out-of-plus-topology} we see
that
$$
R\Gamma(x, H) =
R\Hom_{D(\mathcal{O}(x))}(R\Gamma(x, N), R\Gamma(x, P))
$$
By Lemma \ref{lemma-cohomology-parasitic} we see that $R\Gamma(x, P) = 0$
if $U \to \mathcal{X}$ is flat and hence $R\Gamma(x, H) = 0$ under the same
hypothesis. By Lemma \ref{lemma-cohomology-parasitic-complex} we conclude that
$R\Gamma(\mathcal{X}, H) = 0$ and therefore $\Hom(K, P) = 0$.

\medskip\noindent
Step 3: $\mathcal{A}$ is contained in $Q_\mathcal{X}$.
Let $K$ be an object of $\mathcal{A}$ and let $x \to x'$
be a morphism of $\mathcal{X}_{affine}$. We have to show that
$$
R\Gamma(x', K) \otimes_{\mathcal{O}(x')}^\mathbf{L} \mathcal{O}(x)
\longrightarrow
R\Gamma(x, K)
$$
is a quasi-isomorphism, see footnote in statement of Cohomology on Sites,
Lemma \ref{sites-cohomology-lemma-cartesion-plus-topology}.
By the proof of Lemma \ref{lemma-bousfield-colocalization}
and the discussion in Remark \ref{remark-QCoh-admissible}
we see that $\mathcal{A}$ is the image of the restriction
of $Lg_!$ to $D_\QCoh(\mathcal{O}_{\mathcal{X}_{flat, fppf}})$.
Thus we may assume $K = Lg_!M$ for some $M$ in
$D_\QCoh(\mathcal{O}_{\mathcal{X}_{flat, fppf}})$.
Then the desired equality follow from Lemma \ref{lemma-higher-shriek-QC}.
\end{proof}

\end{document}
